\section{Getting started}

\subsection{Requirements and installation}

\spetc{} needs Python 3.10 or later with Tk support. Install the
dependencies and start the program from the distribution directory:

\begin{verbatim}
python3 -m pip install -r requirements.txt
python3 etc_gui.py
\end{verbatim}

The \code{data/} directory shipped with the program must stay beside
\code{etc\_gui.py}; it contains the stellar template catalogue
(\code{interpola.db.csv} and the CALSPEC/BPGS spectra), the default detector
QE curve (\code{qe.dat}), the SVO filter profiles with their list
(\code{filters.list}), the optional \code{BLANK.xml} unfiltered passband and
the optional Earth-atmosphere FITS table. The data directory can also be
pointed elsewhere from the \textsf{Stellar template selector} panel.

\subsection{First launch}

On start-up \spetc{} restores the last saved configuration from
\code{etc\_user\_config.json} (created automatically; the default preset is
Asiago/Cima Ekar) and the most recently used instrument profile, if any. The
configuration is re-saved after every successful calculation and on normal
exit, so the program always reopens as you left it.

\subsection{A one-minute first calculation}

\begin{enumerate}
  \item Pick a site from the \textsf{Observatory} selector (or type
        latitude, longitude, elevation).
  \item Set the date and UT time, or leave the current time.
  \item Type target coordinates (decimal or sexagesimal), or type a name and
        press \textsf{Resolve name} (needs internet; nothing else ever
        does).
  \item Select an observing filter and a stellar template; enter the target
        magnitude.
  \item Press \textsf{Run ETC}. The Results window opens with the numbers;
        the plot window shows S/N versus time, altitude, and the sky path.
\end{enumerate}
